\documentclass[11pt]{article}
\usepackage[margin=0.85in]{geometry}
\usepackage{amsmath,booktabs,tabularx,microtype,xurl}
\usepackage[hidelinks]{hyperref}
\usepackage{pgfplots}
\pgfplotsset{compat=1.18}
\hypersetup{pdftitle={What Is Missing: A Test for Loss of Sustainable Financing},
  pdfauthor={U.S. Fiscal Dynamics Simulator}}
\setlength{\parindent}{0pt}
\setlength{\parskip}{0.55em}
\title{What Is Missing: A Test for Loss of Sustainable Financing}
\author{U.S. Fiscal Dynamics Simulator}
\date{September 5, 2026}
\begin{document}
\maketitle

\textbf{The missing economic object is a sustainable-financing schedule, together with
a rule for resolving the conflict when the fiscal path cannot be supported.} A debt ledger
shows financing needs. It does not establish the returns at which households and
institutions willingly hold the resulting liabilities, or how expectations of taxes,
inflation, and payments affect those returns.

There is a tractable way to sharpen the question before attempting a large new model.
Mian, Straub, and Sufi show that comparing the current interest rate with growth misses
how borrowing itself changes the rate. Their framework has a maximum sustainable
permanent deficit and a condition stricter than simply interest below growth
\cite{goldilocks}. The diagnostic below is our own annual, single-rate illustration of
this mechanism. It is not a replication or calibration of their model.

\section{The test we should have been running}

Let $b$ be debt/GDP, $d$ the primary deficit as a share of current GDP, $n$ nominal growth,
and $i(b)$ the effective nominal borrowing rate associated with a fully repriced debt stock
of size $b$. Hold the price-stability regime and growth fixed for this experiment. Define
\[
F(b)=\frac{n-i(b)}{1+n}b.
\]
$F(b)$ is the primary deficit compatible with holding that particular debt ratio constant.
The exact annual debt map is
\[
b_{t+1}-b_t=d-F(b_t).
\]
The question becomes whether the fiscal path meets this schedule at a point that attracts
nearby debt paths. That is a different question from whether gross issuance exceeds a
chosen multiple of reference auctions.

With constant rates below growth, $F$ rises linearly and a positive deficit can have a
stable resting point. If rates rise with debt, however, $F$ can reach a maximum and then
fall. More borrowing initially provides additional room for deficit financing; beyond
the maximum, the associated financing costs overwhelm that benefit.

For constant $n$, the slope is
\[
F'(b)=\frac{n-i(b)-b i'(b)}{1+n}.
\]
At a fixed point, the annual transition derivative is $1-F'(b)$. Local stability of this
scalar annual map requires $0<F'(b)<2$. The upper restriction prevents discrete-time
overshooting; with the small coefficients used here it is slack. At the interior maximum,
\[
n-i(b)=b i'(b),
\]
so the borrowing rate can still be below growth when the stable branch disappears.

If $d>\max_{b\geq0}F(b)$, debt increases at every debt ratio. Under these fixed assumptions
there is no bounded debt path. This does not mean a market fails on the first simulated
date, or that the price level must be the variable that adjusts. It means the specified
permanent deficit has no sustainable continuation within this financing environment.

\section{A worked example with no imposed debt ceiling}

Take $n=4\%$ and $i(b)=3\%+0.01(b-1)$. The rate is 3 percent at debt of 100 percent of GDP;
each additional debt/GDP percentage point raises the rate by one basis point. All these
values are illustrative, including the slope. They are not an estimate of U.S. fiscal space.

The resulting schedule is
\[
F(b)=\frac{0.02b-0.01b^2}{1.04}.
\]
It peaks at $b=1$, with a maximum stationary primary deficit of 0.9615 percent of GDP.
A permanent deficit of 0.8 percent intersects it twice: at 59.01-percent debt/GDP, a stable
fixed point, and at 140.99 percent, an unstable one. Starting at 100 percent, the annual
map moves toward the lower fixed point. Starting above 140.99 percent, it moves upward
instead. At the maximum deficit the two roots merge; at a deficit of 1.25 percent there
are no roots. The calculation never imposes a hard debt or auction-volume ceiling.

\begin{figure}[htbp]
\centering
\begin{tikzpicture}
\begin{axis}[
 width=0.94\textwidth,height=0.48\textwidth,
 xmin=0,xmax=220,ymin=-0.25,ymax=1.5,
 xlabel={Debt/GDP, percent},ylabel={Primary deficit/GDP, percent},
 grid=major,grid style={gray!20}
]
\addplot[blue!65!black,very thick] table[col sep=comma,x=debt_gdp_pct,
 y=steady_primary_deficit_pct] {../data/processed/fiscal_space_curve.csv};
\addplot[black,dashed,thick] coordinates {(0,0.8)(220,0.8)};
\addplot[red!75!black,dashed,thick] coordinates {(0,1.25)(220,1.25)};
\addplot[only marks,mark=*,mark size=3pt,green!40!black]
 coordinates {(59.012197,0.8)};
\addplot[only marks,mark=o,mark size=3pt,red!75!black]
 coordinates {(140.987803,0.8)};
\node[font=\small,anchor=west,red!75!black] at (axis cs:4,1.34)
 {1.25\% deficit: no intersection};
\node[font=\small,anchor=south,blue!65!black] at (axis cs:110,1.01)
 {Deficit compatible with constant debt/GDP};
\node[font=\small,anchor=south west] at (axis cs:4,0.81) {0.8\% deficit};
\node[font=\small,anchor=north,green!40!black] at (axis cs:59.012197,0.73) {Stable};
\node[font=\small,anchor=south west,red!75!black] at (axis cs:144,0.81) {Unstable};
\end{axis}
\end{tikzpicture}
\caption{Illustrative financing geometry. Intersections come from the stated rate and
growth assumptions; none is an identified U.S. debt threshold. The high-deficit line has
no intersection even though interest is initially below growth.}
\end{figure}

This is progress in the type of question the code answers. It is not yet progress in
estimating the U.S. threshold: an arbitrary slope can produce an arbitrary maximum.
The shape and location must be disciplined by economic behavior and evidence. The test
also distinguishes a fiscal path with no sustainable resting point from the date at which
investors recognize that problem. Those dates need not coincide.

\section{What would turn the illustration into a U.S. model?}

\textbf{First, explain and estimate the financing schedule.} Investor portfolios must
account for Treasury securities, remunerated reserves, currency, private capital, and
foreign alternatives. Safe-asset and collateral services affect required returns.
Beliefs about future primary balances, inflation, and payment risk affect valuations.
Data on holdings, supply shocks, yields, and returns can constrain this relationship.
Observed Treasury yields alone cannot identify whether buyers expect future fiscal reform
or unusually strong demand for safe assets. Discounting fiscal backing also requires its
risk characteristics, not just today's safe Treasury rate \cite{fiscalCapacity}.

\textbf{Second, connect the schedule to production and the existing cohort ledger.}
Debt can change investment, saving, and growth; taxes and benefits affect private incomes
and the primary balance. Maturity structure determines how marginal rates reach effective
costs. The illustrative annual calculation assumes immediate repricing and cannot measure
that transition. The detailed Treasury ledger is precisely the existing asset to retain.

\textbf{Third, specify which authority changes behavior.} ``Fiscal policy unchanged''
does not specify the Federal Reserve's response to an unfinanceable path. Under a monetary
accommodation branch, solve for prices and real returns consistent with asset demand and
the consolidated budget. Under a price-stability branch, calculate the fiscal adjustment
or payment change needed to restore consistency. Support through interest-bearing reserves
must retain its consolidated financing cost. For each branch, investors' anticipated
response must agree with the response actually simulated. These are alternative conditional
outcomes until regime probabilities have an empirical basis.

\textbf{Fourth, validate what determines the boundary.} Replicate historical responses to
debt supply, monetary tightening, and fiscal news before extrapolating a crisis. Report
how the stable region moves with productivity, saving, foreign demand, and fiscal reaction.
A model that produces collapse for every parameter choice fails to represent sustainable
rollover; a model that always finances every promise at a supplied rate cannot discover
loss of support. The acceptance test is to reproduce both with justified assumptions.

The useful target is therefore an estimate of \emph{which fiscal paths remain supportable
under which economic and policy conditions}, followed by conditional transition paths
when they do not. A calendar crisis forecast is a further task requiring beliefs, shocks,
and policy selection; it cannot be recovered by increasing the number of simulation steps.

\section*{Reproduction}

\begin{verbatim}
uv run python scripts/run_fiscal_space_experiment.py
uv run pytest tests/test_fiscal_space.py
cd docs
latexmk -pdf -interaction=nonstopmode -halt-on-error \
  fiscal_space_mechanism.tex
\end{verbatim}
Parameters are recorded in \path{config/scenarios/fiscal_space_illustration.json};
curves, fixed points, and 80-year illustrative trajectories are
\path{data/processed/fiscal_space_*.csv}. Years are deliberately not mapped to U.S.
calendar dates. The five tests check fixed-point stability, merging and disappearance
of roots, dependence on inherited debt, and failure of the simple interest-below-growth
criterion. They validate the implementation, not an economic calibration.

\begin{thebibliography}{9}
\bibitem{goldilocks}
Atif Mian, Ludwig Straub, and Amir Sufi, ``A Goldilocks Theory of Fiscal Deficits,''
\emph{American Economic Review}, 115(12), 2025, pp. 4253--4291.
\url{https://www.aeaweb.org/articles?id=10.1257/aer.20220308}.
\bibitem{fiscalCapacity}
Zhengyang Jiang, Hanno Lustig, Stijn Van Nieuwerburgh, and Mindy Z. Xiaolan,
\emph{Measuring U.S. Fiscal Capacity using Discounted Cash Flow Analysis},
NBER Working Paper 29902, 2022.
\url{https://www.nber.org/papers/w29902}.
\end{thebibliography}
\end{document}
